\section{El gusto y la no homogeneidad de la inteligencia}

Yang Zhilin (杨植麟) planteó una idea profunda: la siguiente generación de modelos no solo necesita una tecnología más sólida, sino también \textbf{taste} (gusto y sensibilidad estética).

\begin{importantbox}{La inteligencia es un Non-Fungible Token}
La electricidad es fungible (un kilovatio-hora en Shenzhen y uno en Beijing son completamente equivalentes), y el primer centavo de una cuenta bancaria es exactamente igual al último. Pero la inteligencia es \textbf{non-fungible}: la inteligencia de un director general, la de un diseñador y la de un músico son completamente distintas. Construir un modelo consiste, en esencia, en crear una cosmovisión: ¿qué es una buena IA? ¿qué valores debería perseguir?
\end{importantbox}

En el espacio de los agentes, como el test-time scaling genera más tokens, el espacio de diferenciación de la inteligencia crece de forma exponencial: los modelos no convergen hacia lo mismo, sino que surgen cada vez más taste nuevos.

\begin{knowledgebox}{El eco de Jobs}
Yang Zhilin (杨植麟) citó la célebre frase de Jobs, ``It all comes down to taste'', y sostuvo que se aplica igual al campo de la IA. En la dimensión de la inteligencia existe un espacio de gusto inmenso, y es precisamente ahí donde cada equipo de modelos puede construir una ventaja competitiva propia.
\end{knowledgebox}

